\chapter{Zaawansowana Współbieżność i pobieranie danych (Wykład 3)}

\section{Ustrukturyzowana Współbieżność}

Poprzednio nauczyliśmy się, jak uruchamiać zadania w tle (\textit{nastawiać wodę}) i nie blokować przy tym kuchni. Co w sytuacji, gdy przygotowanie dania wymaga wykonania \textbf{kilku czynności na raz}, a my musimy mieć pewność, że wszystkie zostały ukończone, zanim wydamy posiłek?

Tu z pomocą przychodzi funkcja \texttt{join()}, która jest fundamentem synchronizacji w świecie korutyn.

Wyobraź sobie, że jesteś Szefem Kuchni (Korutyna-Rodzic). Masz przygotować główne danie: \textit{Stek z warzywami}. Nie będziesz robić wszystkiego sam.
\begin{itemize}
\item Wołasz Pomocnika nr 1: \textit{Usmaż mięso!} (zajmie to 2 sekundy).
\item Wołasz Pomocnika nr 2: \textit{Ugotuj warzywa!} (zajmie to 3 sekundy).
\end{itemize}

Obaj pomocnicy ruszają do pracy w tym samym momencie (współbieżność). Główna korutyn jest wolna i może zająć się innymi sprawami. Ale nie może wydać dania dopóki nie będzie gotowe. To \textit{czekanie} wykonujemy za pomocą \texttt{join()}.

Zobaczmy kod, któy symuluje taką sytuację:

\begin{minted}[frame=single,framesep=10pt]{kotlin}
// 1. Szef Kuchni (Rodzic) rozpoczyna pracę
scope.launch {
    logs.add("Szef kuchni (rodzic): Zaczynamy!")

    // 2. Zlecenie zadania Pomocnikowi 1 (Dziecko 1)
    val jobMieso = launch {
        delay(2000) // Symulacja smażenia
        logs.add("Kucharz 1: Mięso usmażone (2s).")
    }

    // 3. Zlecenie zadania Pomocnikowi 2 (Dziecko 2)
    val jobWarzywa = launch {
        delay(3000) // Symulacja gotowania
        logs.add("Kucharz 2: Warzywa gotowe (3s).")
    }

    logs.add("Szef kuchni: Zadania zlecone, scope czeka na zakończenie...")

    // 4. synchronizacja: Szef czeka na wyniki
    jobWarzywa.join()
    jobMieso.join()

    // 5. Finał
    logs.add("Szef kuchni: Wszyscy skończyli! Można podawać danie.")
}
\end{minted}

W przykładzie widzimy koordynację zadań podrzędnych. Cały proces rozpoczyna się w momencie, gdy korutyna nadrzędna - nasz metaforyczny Szef Kuchni - uruchamia podzadania za pomocą funkcji \texttt{launch}. Jest to swego rodzaju \textit{bilet zamówienia}, który daje nam unikalny uchwyt do konkretnej, trwającej korutyny. W naszym przypadku tworzymy dwa takie uchwyty: \texttt{jobMieso} oraz \texttt{jobWarzywa}, które stają się dziećmi głównej korutyny uruchomionej w scope.

Oba zadania startują niemal w tym samym momencie, co oznacza, że nie czekamy z gotowaniem warzyw do momentu, aż mięso będzie gotowe. Dzięki temu całkowity czas operacji nie jest sumą czasów poszczególnych zadań (5 sekund), lecz odpowiada czasowi najdłuższego z nich (3 sekundy). Jednak ta niezależność podzadań rodzi problem \textbf{synchronizacji}: co, jeśli Szef Kuchni zakończy swoją pracę szybciej niż jego pomocnicy? Bez odpowiedniej kontroli, log \textit{Można podawać danie} pojawiłby się natychmiast po zleceniu zadań, co w naszej metaforze oznaczałoby wydanie klientowi pustego talerza, zanim składniki zdążą się ugotować.

Rozwiązaniem problemu wyścigu (race condition) jest funkcja \texttt{join()}. Z technicznego punktu widzenia jest to funkcja zawieszająca (\texttt{suspend function}), która służy do synchronizacji cyklu życia korutyn. Kiedy korutyna nadrzędna (rodzic) napotyka instrukcję \texttt{jobWarzywa.join()}, jej wykonanie zostaje zawieszone. Należy wyraźnie podkreślić różnicę między zawieszeniem a zablokowaniem: wątek obsługujący tę korutynę nie jest blokowany i może w tym czasie wykonywać inne operacje systemowe. Mechanizm \texttt{join} wprowadza korutynę wywołującą w stan oczekiwania, który trwa dopóki obserwowany obiekt \texttt{Job} nie osiągnie stanu końcowego (Completed lub Cancelled). Dopiero gdy podzadanie faktycznie zakończy swoje działanie, maszyna stanów wznawia korutynę rodzica od kolejnej linii kodu. Dzięki jawnemu wywołaniu \texttt{join()} na obu obiektach \texttt{Job}, realizujemy kontrakt Ustrukturyzowanej Współbieżności, w którym rodzic świadomie koordynuje pracę swoich dzieci i gwarantuje, że żadna operacja nie zostanie zakończona przedwcześnie, zapewniając spójność danych i przewidywalność przepływu aplikacji.

Warto zatrzymać się przy aspekcie, który czyni ustrukturyzowaną współbieżność \textit{bezpieczną} - jest to \textbf{automatyczna propagacja anulowania} - wpominaliśmy już o tym wcześniej. Ale omówmy to jeszcze raz.

Wyobraźmy sobie, że nasz Szef Kuchni (rodzic) otrzymuje informację, że klient anulował zamówienie i wychodzi z restauracji. W świecie \textbf{ustrukturyzowanej współbieżności} Szef nie musi biegać po kuchni i osobiście prosić każdego pomocnika o przerwanie pracy. Wywołanie \texttt{job.cancel()} na korutynie-rodzicu \textbf{automatycznie} przesyła sygnał anulowania do wszystkich dzieci (\texttt{jobMieso}, \texttt{jobWarzywa}). Jeśli pomocnicy wykonują funkcje zawieszające (takie jak \texttt{delay}), natychmiast przerwą pracę.

W programowaniu aplikacji na Androida ten mechanizm jest niezwykle istotny. Kiedy użytkownik zamyka ekran (\texttt{Activity}), powiązany z nim zakres (\texttt{viewModelScope} lub \texttt{lifecycleScope}) zostaje anulowany. Dzięki hierarchii rodzic-dziecko, wszystkie trwające zapytania sieciowe czy obliczenia w tle zostają \textit{posprzątane} automatycznie. Zapobiega to wyciekom pamięci i marnowaniu baterii na procesy, których wyniku nikt już nie zobaczy.

\section{Odbieranie Wyników: Mechanizm \texttt{async} i \texttt{await}}

Mimo że funkcja \texttt{launch} jest niezwykle przydatna, posiada jedno istotne ograniczenie: działa na zasadzie \textit{odpal i zapomnij} (\textit{fire-and-forget}). Zwraca ona obiekt \texttt{Job}, który pozwala nam sprawdzić czy praca trwa, lub ją przerwać, ale nie pozwala nam \textit{wyjąć} ze środka żadnej wartości. W świecie rzeczywistym nie zawsze chcemy tylko zlecić zadanie; często potrzebujemy, aby pracownik coś nam przyniósł (np. konkretny składnik z magazynu). Do takich zadań służy funkcja \texttt{async}.

\subsection{Analogia: Kurier i Pusta Skrzynka} 
Zanim przeanalizujemy kod, posłużmy się analogią zamówienia przesyłki: 
\begin{itemize} 
    \item \textbf{async:} To moment kliknięcia \textit{Kup teraz}. Sklep nie daje nam towaru natychmiast, ale daje nam \textbf{numer śledzenia przesyłki} (to jest nasz obiekt \texttt{Deferred}). 
    \item \textbf{Deferred:} To metaforyczne \textit{puste pudełko}. Pudełko istnieje fizycznie (mamy do niego referencję w kodzie), ale w tej chwili nie możemy z niego skorzystać, bo zawartość jest jeszcze w drodze. 
    \item \textbf{await():} To moment, w którym czekamy na dzwonek kuriera. Jeśli kurier już był - odbieramy zawartość natychmiast. Jeśli jeszcze jedzie - zawieszamy nasze inne czynności i czekamy, aż pudełko zostanie \textit{odpakowane}. 
\end{itemize}

\subsection{Przykład Praktyczny: Generator Tajemniczego Hasła} 
Rozważmy implementację ekranu, który symuluje złożony proces generowania bezpiecznego hasła. Proces ten trwa 3 sekundy i musi zwrócić wynik bezpośrednio do zmiennej stanu interfejsu.

\begin{minted}[frame=single,framesep=10pt]{kotlin}
@Composable fun PasswordGeneratorScreen() { 
    var password by remember { mutableStateOf("...") } 
    val scope = rememberCoroutineScope()

    // Funkcja symulująca ciężkie obliczenia w tle
    suspend fun generatePassword(): String {
        delay(3000) // Symulacja pracy
        return "Kotlin-Is-Super-Secure-123"
    }

    Column(
        modifier = Modifier.fillMaxSize().padding(16.dp),
        horizontalAlignment = Alignment.CenterHorizontally,
        verticalArrangement = Arrangement.Center
    ) {
        Text("Generator Tajemniczego Hasła", style = MaterialTheme.typography.headlineSmall)
        Spacer(Modifier.height(20.dp))

        Button(onClick = {
            password = "Generowanie..."
            
            // Uruchamiamy korutynę nadrzędną (szkielet operacji)
            scope.launch {
                // 1. Zlecenie zadania (async) - natychmiast dostajemy 'pudełko' Deferred
                val deferredPassword: Deferred<String> = async { generatePassword() }
                
                // 2. Oczekiwanie (await) - korutyna zawiesza się tutaj na 3 sekundy
                val result = deferredPassword.await()
                
                // 3. Po 3 sekundach korutyna wznawia pracę z gotowym wynikiem
                password = result
            }
        }) {
            Text("Wygeneruj hasło")
        }

        Spacer(Modifier.height(20.dp))
        Text(password, fontSize = 20.sp, fontFamily = FontFamily.Monospace)
    }
} 
\end{minted}

W powyższym kodzie kluczowy jest moment kliknięcia przycisku. Wywołujemy \path{scope.launch}, aby wejść do świata współbieżności. Wewnątrz niego dzieją się dwie rzeczy: Po pierwsze, linia \texttt{async { generatePassword() }} nie blokuje wykonania kodu. Funkcja \texttt{async} natychmiast zwraca obiekt typu \texttt{Deferred<String>}. Gdybyśmy w tym momencie sprawdzili wartość \texttt{deferredPassword}, dowiedzielibyśmy się jedynie, że zadanie jest w toku. Po drugie, wywołanie \texttt{deferredPassword.await()} jest instrukcją dla mechanizmu korutyn: \textit{Zatrzymaj tę konkretną korutynę w tym miejscu i czekaj, aż} \texttt{deferredPassword} \textit{dostarczy wartość}. Co niezwykle ważne - \textbf{wątek UI pozostaje wolny}. Użytkownik widzi napis \texttt{Generowanie...}, ale aplikacja nie jest zamrożona; system może w tym czasie obsłużyć inne zdarzenia.

Gdy funkcja \texttt{generatePassword()} skończy pracę i zwróci \texttt{String}, \texttt{await()} \textit{odpakuje} go i przypisuje do zmiennej \texttt{result}. Dopiero wtedy wykonuje się ostatnia linia: \texttt{password = result}, co powoduje rekompozycję interfejsu i wyświetlenie hasła.

\section{Most między światami: runBlocking}

Wszystkie funkcje, które poznaliśmy do tej pory - \texttt{launch} oraz \texttt{async} - wymagają do działania istniejącego zakresu (\texttt{CoroutineScope}). Istnieją jednak sytuacje, w których musimy poczekać, aż świat korutyn zakończy swoją pracę, zanim pozwolimy światu zewnętrznemu (blokującemu) ruszyć z miejsca. Do tego celu służy \texttt{runBlocking}. Jak sama nazwa wskazuje, jest to funkcja, która \textbf{blokuje bieżący wątek}, na którym została wywołana, dopóki wszystkie korutyny wewnątrz jej bloku nie zostaną zakończone.

Wyobraź sobie strumień kodu asynchronicznego jako płynącą rzekę - korutyny płyną swobodnie, nie przeszkadzając sobie nawzajem. \texttt{runBlocking} jest jak \textbf{tama}, którą stawiamy w poprzek rzeki. 
\begin{itemize} 
    \item Woda (wątek) przestaje płynąć dalej za tamę. 
    \item Cały nurt zostaje zatrzymany tak długo, aż praca wewnątrz tamy nie zostanie ukończona. 
    \item Dopiero gdy ostatnia korutyna w bloku \texttt{runBlocking} skończy działanie, tama zostaje otwarta i wątek może kontynuować wykonywanie kolejnych linii kodu pod blokiem. 
\end{itemize}

\subsection{Zastosowanie} 
Ponieważ \texttt{runBlocking} blokuje wątek, wywołanie go na wątku UI spowoduje natychmiastowe zamrożenie interfejsu użytkownika na cały czas trwania operacji wewnątrz bloku. Jest to powrót do problemu, który korutyny miały rozwiązać.

Kiedy zatem jest to przydatne? 
\begin{enumerate} 
    \item \textbf{Testy jednostkowe:} W testach chcemy, aby proces testowy poczekał na wynik operacji asynchronicznej, zanim sprawdzi poprawność wyniku (asercję). 
    \item \textbf{Funkcja main():} W prostych programach konsolowych Kotlin, gdzie musimy \textit{zatrzymać} program przed zamknięciem, dopóki korutyny nie skończą pracy. 
    \item \textbf{Integracja z kodem blokującym:} Gdy musimy wywołać kod asynchroniczny wewnątrz biblioteki, która nie wspiera korutyn i wymaga natychmiastowego zwrotu wartości. 
\end{enumerate}

\subsection{Porównanie mechanizmów} 
Poniższy przykład pokazuje, jak \texttt{runBlocking} wpływa na przepływ programu w porównaniu do \texttt{launch}.

\begin{minted}[frame=single,framesep=10pt]{kotlin} 
fun runBlockingExample() { 
    println("1. Start programu")

    // Ten blok zatrzyma wątek na 2 sekundy!
    runBlocking {
        println("2. Wewnątrz runBlocking - start")
        delay(2000)
        println("3. Wewnątrz runBlocking - koniec")
    }

    println("4. Koniec programu - ta linia czekała na tamę")
} 
\end{minted}

W powyższym kodzie napis \textit{4. Koniec programu} pojawi się dopiero po 2 sekundach. Gdybyśmy zamiast \texttt{runBlocking} użyli \texttt{launch} (w odpowiednim \texttt{scope}), napis \textit{4} pojawiłby się natychmiast po napisie \textit{2}, ponieważ rzeka płynęłaby dalej, podczas gdy korutyna pracowałaby \textit{obok}.